% =====================================================================
% RASCUNHO -- seção "Setup de simulação" reorganizada por módulo do simulador
% ---------------------------------------------------------------------
% Incluído em 9-results.tex por \input{chapters/rascunho_setup_simulacao}.
% Substitui o setup antigo (cópia em rascunho_setup_antigo.tex) e
% rascunho_setup_tabelas.tex (as tabelas estão copiadas aqui).
% A entrada tomasin-interference precisa ir para o biblio.bib.
%
% Marcas:
%   % AUTORA (9-results.tex:NNN)  -> texto seu, sem alteração, só mudou de lugar
%   % NOVO                        -> texto proposto; cada frase tem FONTE ou SEM FONTE
%
% Chaves de citação: zhao2025jsac (já existe); tomasin-interference (em
% constraints/rascunho_referencias.bib). ATENÇÃO: rascunho_bloco_canal.tex
% usa a chave antiga zhao-lmmse; trocar por zhao2025jsac se for aproveitado.
% "Lnnn" = linha do paper.md do artigo base.
% =====================================================================

\section{Setup de simulação}

% NOVO -- substitui 9-results.tex:4
Os parâmetros da simulação estão organizados de acordo com a parte do simulador em que são usados: o canal, o sinal e o ruído, o receptor e a avaliação. Os parâmetros do sistema seguem o modelo de \citeonline{zhao2025jsac}, que propõe os equalizadores descentralizados avaliados neste trabalho. Em cada tabela, a coluna de origem indica a referência de cada valor.
% TOM: Zhao et al. é a referência do modelo, não objeto de reprodução.


% ---------------------------------------------------------------------
\subsection{Configurações simuladas}
% ---------------------------------------------------------------------
% Vem primeiro: M, C, K, r, IoT e modulação mudam entre as configurações e
% não se repetem nas tabelas seguintes.

% AUTORA (9-results.tex:214), tom ajustado
Foram simuladas seis configurações, as mesmas avaliadas por \citeonline{zhao2025jsac}. A configuração 1 é a de referência. Cada uma das outras muda um parâmetro em relação a ela. A Tabela~\ref{tab:paineis} mostra essas configurações.

\begin{table}[htbp]
\color{red}
\caption{Configurações simuladas.}
\label{tab:paineis}
\centering
\footnotesize
\renewcommand{\arraystretch}{1.2}
\begin{tabular}{|c|c|c|c|c|c|c|}
\hline
Configuração & $M$ & $C$ & $M_c$ & $K$ & IoT (dB) & Modulação \\
\hline
-- & \texttt{M} & \texttt{C} & -- & \texttt{K} & \texttt{iot\_db} & \texttt{modulation} \\
\hline
1 & 128 & 8  & 16 & 8  & 10 & 16-QAM \\
\hline
2 & 256 & 8  & 32 & 8  & 10 & 16-QAM \\
\hline
3 & 256 & 16 & 16 & 8  & 10 & 16-QAM \\
\hline
4 & 128 & 8  & 16 & 12 & 10 & 16-QAM \\
\hline
5 & 128 & 8  & 16 & 8  & 20 & 16-QAM \\
\hline
6 & 128 & 8  & 16 & 8  & 10 & QPSK \\
\hline
\end{tabular}
\legend{\textcolor{red}{Fonte: adaptado de \citeonline{zhao2025jsac}. Em todas as configurações, $r = 8$. Nomes no código: campos de \texttt{Panel} (\texttt{dbp\_sim/params.py}); $M_c = M/C$.}}
\end{table}

% AUTORA (9-results.tex:245) -- parágrafo retirado do texto (tom); o conteúdo ficou em comentário.
% Configurações 1 a 6 = painéis (a), (b), (c), (e), (f), (g) da Fig. 8 de Zhao et al.
% O painel (d) não aparece no artigo. Na configuração 4 (painel e), r = 8 foi
% mantido; L577 liga r a K, então r = 12 também seria defensável.


% ---------------------------------------------------------------------
\subsection{Canal}\label{subsec:setup_canal}
% ---------------------------------------------------------------------

% NOVO
A Tabela~\ref{tab:param_canal} mostra os parâmetros usados para gerar as matrizes de canal $\mathbf{H}$ e $\overline{\mathbf{H}}$.
% AUTORA (9-results.tex:67), adaptado: "A Tabela ... lista esses valores" saiu,
% porque agora os padrões estão na mesma tabela.
Os demais parâmetros do QuaDRiGa foram mantidos no valor padrão do simulador \cite{quadriga-manual}.
% Saiu "e não mudam o resultado de forma importante" (não foi medido).

\begin{table}[htbp]
\caption{Parâmetros do cenário e do canal.}
\label{tab:param_canal}
\centering
\footnotesize
\renewcommand{\arraystretch}{1.0} % 1.0 (e não 1.2) para a legenda caber na página
\begin{tabular}{|p{2.2cm}|p{3.6cm}|p{4.6cm}|p{3.4cm}|}
\hline
\raggedright Origem & \raggedright Parâmetro & \raggedright Nome no código & \raggedright Valor \tabularnewline
\hline
\raggedright \cite{zhao2025jsac} & \raggedright Número de células e de BSs & \raggedright \texttt{M2\_\allowbreak{}Channel.\allowbreak n\_\allowbreak{}cells} & \raggedright 1 \tabularnewline
\cline{2-4}
 & \raggedright Setor onde ficam os UEs & \raggedright \texttt{M2\_\allowbreak{}Channel.\allowbreak sector\_\allowbreak{}deg} & \raggedright $120^\circ$ \tabularnewline
\cline{2-4}
 & \raggedright Distância dos UEs até a BS & \raggedright \texttt{M2\_\allowbreak{}Channel.\allowbreak radius\_\allowbreak{}min\_\allowbreak{}m}, \texttt{radius\_\allowbreak{}max\_\allowbreak{}m} & \raggedright 50 a 100 m \tabularnewline
\cline{2-4}
 & \raggedright Posição dos interferentes & \raggedright \texttt{M2\_\allowbreak{}Channel.\allowbreak interferers\_\allowbreak{}like\_\allowbreak{}targets} & \raggedright Mesma regra dos UEs atendidos \tabularnewline
\cline{2-4}
 & \raggedright Gerador do canal & \raggedright -- & \raggedright QuaDRiGa \tabularnewline
\cline{2-4}
 & \raggedright Canal constante dentro de um PRB & \raggedright -- & \raggedright Canal plano \tabularnewline
\hline
\raggedright Padrão do QuaDRiGa \cite{quadriga-manual} & \raggedright Usa só a versão mínima do modelo 3GPP & \raggedright \texttt{qd\_\allowbreak{}simulation\_\allowbreak{}parameters.\allowbreak use\_\allowbreak{}3GPP\_\allowbreak{}baseline} & \raggedright \texttt{false} \tabularnewline
\cline{2-4}
 & \raggedright Forma da correlação espacial dos parâmetros de larga escala & \raggedright \texttt{qd\_\allowbreak{}simulation\_\allowbreak{}parameters.\allowbreak autocorrelation\_\allowbreak{}function} & \raggedright \texttt{Comb300} \tabularnewline
\cline{2-4}
 & \raggedright Soma o atraso do caminho direto aos atrasos & \raggedright \texttt{qd\_\allowbreak{}simulation\_\allowbreak{}parameters.\allowbreak use\_\allowbreak{}absolute\_\allowbreak{}delays} & \raggedright \texttt{false} \tabularnewline
\cline{2-4}
 & \raggedright Cada caminho começa com uma fase aleatória & \raggedright \texttt{qd\_\allowbreak{}simulation\_\allowbreak{}parameters.\allowbreak use\_\allowbreak{}random\_\allowbreak{}initial\_\allowbreak{}phase} & \raggedright \texttt{true} \tabularnewline
\cline{2-4}
 & \raggedright Amostras por meio comprimento de onda na trajetória & \raggedright \texttt{qd\_\allowbreak{}simulation\_\allowbreak{}parameters.\allowbreak sample\_\allowbreak{}density} & \raggedright 2,5 \tabularnewline
\cline{2-4}
 & \raggedright Altura da BS & \raggedright \texttt{qd\_\allowbreak{}layout.\allowbreak tx\_\allowbreak{}position} & \raggedright 25 m \tabularnewline
\hline
\raggedright Este trabalho & \raggedright Forma de uso do QuaDRiGa & \raggedright \texttt{M2\_\allowbreak{}Channel.\allowbreak backend} & \raggedright Canais exportados em arquivo e lidos em Python \tabularnewline
\cline{2-4}
 & \raggedright Versão do QuaDRiGa & \raggedright \texttt{M2\_\allowbreak{}Channel.\allowbreak quadriga\_\allowbreak{}version} & \raggedright 2.8.1-0 \tabularnewline
\cline{2-4}
 & \raggedright Cenário & \raggedright \texttt{M2\_\allowbreak{}Channel.\allowbreak quadriga\_\allowbreak{}scenario} & \raggedright \texttt{3GPP\_\allowbreak{}38.901\_\allowbreak{}UMa\_\allowbreak{}NLOS} \tabularnewline
\cline{2-4}
 & \raggedright Frequência da portadora & \raggedright \texttt{M2\_\allowbreak{}Channel.\allowbreak center\_\allowbreak{}frequency\_\allowbreak{}hz} & \raggedright 3,5 GHz \tabularnewline
\cline{2-4}
 & \raggedright Largura de banda & \raggedright \texttt{M2\_\allowbreak{}Channel.\allowbreak bandwidth\_\allowbreak{}hz} & \raggedright 20 MHz \tabularnewline
\cline{2-4}
 & \raggedright Extração do canal plano & \raggedright \texttt{M2\_\allowbreak{}Channel.\allowbreak flat\_\allowbreak{}channel\_\allowbreak{}extraction} & \raggedright Resposta em uma subportadora \tabularnewline
\cline{2-4}
 & \raggedright Arranjo da BS & \raggedright \texttt{M2\_\allowbreak{}Channel.\allowbreak array\_\allowbreak{}geometry} & \raggedright ULA com espaçamento $\lambda/2$ \tabularnewline
\cline{2-4}
 & \raggedright Elementos de antena (BS e UE) & \raggedright \texttt{M2\_\allowbreak{}Channel.\allowbreak array\_\allowbreak{}geometry} & \raggedright Omnidirecionais \tabularnewline
\cline{2-4}
 & \raggedright Orientação do setor & \raggedright \texttt{M2\_\allowbreak{}Channel.\allowbreak array\_\allowbreak{}geometry} & \raggedright Centrado na direção perpendicular ao arranjo \tabularnewline
\cline{2-4}
 & \raggedright Altura do UE & \raggedright \texttt{M2\_\allowbreak{}Channel.\allowbreak ue\_\allowbreak{}height\_\allowbreak{}m} & \raggedright 1,5 m \tabularnewline
\cline{2-4}
 & \raggedright Posição dos UEs & \raggedright \texttt{M2\_\allowbreak{}Channel.\allowbreak ue\_\allowbreak{}placement} & \raggedright Uniforme na área do setor \tabularnewline
\cline{2-4}
 & \raggedright Realizações & \raggedright \texttt{M2\_\allowbreak{}Channel.\allowbreak independent\_\allowbreak{}drops} & \raggedright Posições novas a cada realização \tabularnewline
\cline{2-4}
 & \raggedright Normalização do canal & \raggedright \texttt{M2\_\allowbreak{}Channel.\allowbreak normalization} & \raggedright Controle de potência por UE \tabularnewline
\hline
\end{tabular}
\legend{Fonte: \citeonline{zhao2025jsac}, código-fonte do QuaDRiGa v2.8.1 e a autora. $M$, $K$ e $r$ estão na Tabela~\ref{tab:paineis}. Nomes no código: \texttt{dbp\_sim/params.py} e propriedades das classes do QuaDRiGa.} % TODO citar QuaDRiGa
\end{table}
% FONTES (artigo): L577 (célula única, setor, raio, QuaDRiGa, interferentes
%   "molded similarly"); L614 (canal constante no PRB).
% Altura do UE: o script fixa 1,5 m, igual ao padrão de randomize_rx_positions;
%   por isso está em "Escolhido" e não em "Padrão do QuaDRiGa".
% Altura da BS: o script também fixa 25 m, igual ao padrão de qd_layout.
%   Mantida em "Padrão" porque a escolha do cenário se apoia nesse padrão.
%   Se preferir coerência com a altura do UE, mover para "Escolhido".

% AUTORA (9-results.tex:173-175) -- antes na subseção "Cenário escolhido".
% A frase da linha 171 ("O QuaDRiGa define cada cenário por um arquivo de
% configuração...") saiu: o mesmo conceito já está na fundamentação
% (0.channel-modelling.tex, "Cenário de propagação").
Neste trabalho foi usado o cenário \texttt{3GPP\_38.901\_UMa\_NLOS}, que representa uma macrocélula urbana (\textit{Urban Macro} -- UMa) sem linha de visada (NLoS) entre a BS e os UEs. % TODO citar 3GPP TR 38.901 e QuaDRiGa
Os cenários UMa e UMi (\textit{Urban Micro}) aceitam a mesma faixa de distâncias, de 10 a 5000 m. Então o raio de 50 a 100 m não separa os dois. A diferença está na altura da BS: 10 m no UMi e 25 m no UMa. Foi usada a altura padrão do QuaDRiGa, de 25 m, que corresponde ao cenário UMa. % TODO citar QuaDRiGa (qd_layout: tx em [0;0;25])
A Tabela~\ref{tab:cenario_uma} compara a faixa de validade do cenário com os valores usados na simulação.

% AUTORA (9-results.tex:178-209) -- tabela de validade do cenário
\begin{table}[htbp]
\color{red}
\caption{Faixa de validade do cenário \texttt{3GPP\_38.901\_UMa\_NLOS} e valores usados na simulação.}
\label{tab:cenario_uma}
\centering
\footnotesize
\renewcommand{\arraystretch}{1.2}
\begin{tabular}{|p{3.2cm}|p{4.6cm}|p{3.1cm}|p{2.3cm}|}
\hline
\raggedright Parâmetro & \raggedright Nome no código & \raggedright Faixa válida do cenário & \raggedright Valor usado \tabularnewline
\hline
\raggedright Frequência da portadora & \raggedright \texttt{M2\_\allowbreak{}Channel.\allowbreak center\_\allowbreak{}frequency\_\allowbreak{}hz} & \raggedright 0,5 a 100 GHz & \raggedright 3,5 GHz \tabularnewline
\hline
\raggedright Distância BS--UE & \raggedright \texttt{M2\_\allowbreak{}Channel.\allowbreak radius\_\allowbreak{}min\_\allowbreak{}m}, \texttt{radius\_\allowbreak{}max\_\allowbreak{}m} & \raggedright 10 a 5000 m & \raggedright 50 a 100 m \tabularnewline
\hline
\raggedright Altura da BS & \raggedright \texttt{M2\_\allowbreak{}Channel.\allowbreak bs\_\allowbreak{}height\_\allowbreak{}m} & \raggedright 25 m & \raggedright 25 m \tabularnewline
\hline
\raggedright Altura do UE & \raggedright \texttt{M2\_\allowbreak{}Channel.\allowbreak ue\_\allowbreak{}height\_\allowbreak{}m} & \raggedright 1,5 a 22,5 m & \raggedright 1,5 m \tabularnewline
\hline
\raggedright Número de \textit{clusters} & \raggedright \texttt{NumClusters} & \multicolumn{2}{c|}{21} \tabularnewline
\hline
\raggedright Caminhos por \textit{cluster} & \raggedright \texttt{NumSubPaths} & \multicolumn{2}{c|}{20} \tabularnewline
\hline
\raggedright Desvio padrão do sombreamento & \raggedright \texttt{SF\_\allowbreak{}sigma} & \multicolumn{2}{c|}{6 dB} \tabularnewline
\hline
\end{tabular}
\legend{\textcolor{red}{Fonte: arquivo \texttt{3GPP\_38.901\_UMa\_NLOS.conf} do QuaDRiGa v2.8.1. Nomes no código: \texttt{dbp\_sim/params.py} e chaves do arquivo \texttt{.conf}.}} % TODO citar QuaDRiGa
\end{table}





% AUTORA (9-results.tex:153)
A frequência da portadora não é definida pelo arquivo do cenário. Ela precisa ser escolhida e informada, porque os parâmetros do cenário mudam com a frequência. Foi usado 3,5 GHz. % TODO: 3,5 GHz não tem fonte. Sugestão: citar uma referência que use 3,5 GHz como banda típica do 5G NR (ex.: livro 5G and Beyond), ou dizer só que é uma escolha deste trabalho.

% AUTORA (9-results.tex:155-156)
O modelo de sistema considera o canal plano, ou seja, o canal é o mesmo em todas as subportadoras de um PRB \cite{zhao2025jsac}. % L614
O QuaDRiGa gera a resposta do canal em frequência. Por isso foi usada a resposta em uma única subportadora como a matriz $\mathbf{H}$. A largura de banda de 20 MHz só define a grade de subportadoras e não muda o resultado.

% AUTORA (9-results.tex:158)
Foi usado um arranjo linear uniforme (ULA) de $M$ elementos, com espaçamento de meio comprimento de onda ($\lambda/2$). Todos os elementos são omnidirecionais, ou seja, recebem com o mesmo ganho em todas as direções. Os UEs também têm uma antena omnidirecional.

% AUTORA (9-results.tex:160-161)
Os UEs são sorteados de forma uniforme na área do setor de $120^\circ$, entre 50 e 100 m da BS. O centro do setor fica na direção perpendicular ao arranjo. O ângulo é sorteado de forma uniforme dentro do setor. A distância é sorteada de forma que a densidade de UEs seja a mesma em toda a área. Assim a região mais distante, que tem área maior, recebe mais UEs. Os interferentes seguem a mesma regra dos UEs atendidos \cite{zhao2025jsac}. % L577
Cada uma das 100 realizações sorteia posições novas para todos os UEs.

% AUTORA (9-results.tex:163-164)
Como os UEs ficam entre 50 e 100 m da BS, a perda de percurso varia muito entre eles. Nos canais gerados, a diferença de potência entre o UE mais forte e o mais fraco teve mediana de 13 dB nas 100 realizações. Sem nenhum ajuste, a taxa de erro média fica dominada pelo UE mais fraco. Por isso cada coluna de $\mathbf{H}$ e $\overline{\mathbf{H}}$ foi normalizada para ter a mesma potência média, com $\mathbb{E}[|h_{mk}|^2] = 1$. Isso equivale a um controle de potência ideal no enlace de subida.
% Medição deste trabalho (13 dB). Saíram do texto: "foi o parâmetro que mais mudou o
%   resultado", "o artigo não fala de controle de potência" e as razões de SER contra a
%   Fig. 8(a) (0,56-0,65x / 0,8-7,5x). Se forem usadas, o lugar é a discussão dos
%   resultados, não o setup. Fonte possível para controle de potência: bjornson-mimo,
%   p. 311 (ver rascunho_bloco_canal.tex).


% ---------------------------------------------------------------------
\subsection{Sinal e ruído}\label{subsec:setup_ruido}
% ---------------------------------------------------------------------

% NOVO
A Tabela~\ref{tab:param_ruido} mostra os parâmetros usados para gerar os símbolos transmitidos e o ruído.

\begin{table}[htbp]
\caption{Parâmetros do sinal e do ruído.}
\label{tab:param_ruido}
\centering
\footnotesize
\renewcommand{\arraystretch}{1.2}
\begin{tabular}{|p{2.2cm}|p{3.6cm}|p{4.6cm}|p{3.4cm}|}
\hline
\raggedright Origem & \raggedright Parâmetro & \raggedright Nome no código & \raggedright Valor \tabularnewline
\hline
\raggedright \cite{zhao2025jsac} & \raggedright Modelo do ruído & \raggedright \texttt{M5\_\allowbreak{}Noise.\allowbreak noise\_\allowbreak{}model} & \raggedright $\mathbf{n} \sim \mathcal{CN}(\mathbf{0}, \beta\overline{\mathbf{H}}\,\overline{\mathbf{H}}^H + \sigma^2\mathbf{I})$ \tabularnewline
\cline{2-4}
 & \raggedright Definição da SNR & \raggedright \texttt{M5\_\allowbreak{}Noise.\allowbreak snr\_\allowbreak{}definition} & \raggedright $10\log_{10}(E_s/\sigma^2)$ \tabularnewline
\cline{2-4}
 & \raggedright Energia média do símbolo & \raggedright \texttt{M4\_\allowbreak{}Source.\allowbreak Es} & \raggedright $E_s = 1$ \tabularnewline
\cline{2-4}
 & \raggedright Potência do ruído térmico & \raggedright \texttt{M5\_\allowbreak{}Noise.\allowbreak snr\_\allowbreak{}definition} & \raggedright $\sigma^2 = 10^{-\mathrm{SNR}/10}$ \tabularnewline
\hline
\raggedright Este trabalho & \raggedright Energia da constelação & \raggedright \texttt{M4\_\allowbreak{}Source.\allowbreak unit\_\allowbreak{}energy\_\allowbreak{}constellation} & \raggedright Média unitária \tabularnewline
\cline{2-4}
 & \raggedright Potência dos interferentes & \raggedright \texttt{M5\_\allowbreak{}Noise.\allowbreak beta\_\allowbreak{}convention} & \raggedright $\beta = \sigma^2(10^{\mathrm{IoT}/10} - 1)/r$ \tabularnewline
\cline{2-4}
 & \raggedright Ruído dos símbolos de dados & \raggedright \texttt{M5\_\allowbreak{}Noise.\allowbreak data\_\allowbreak{}noise\_\allowbreak{}independent} & \raggedright Independente das $N$ amostras \tabularnewline
\hline
\end{tabular}
\legend{Fonte: \citeonline{zhao2025jsac} e a autora. IoT e modulação estão na Tabela~\ref{tab:paineis}. Nomes no código: \texttt{dbp\_sim/params.py}.}
\end{table}
% FONTES: modelo do ruído L113, L577; SNR L577; E_s e sigma^2 L577, L79.
% beta: IoT definido em L577, N_0 nunca definido, lido como sigma^2 via [P] L207,
%   e dividido por r para que a IoT medida por antena seja a IoT pedida.
%   A fórmula com /r vale com a normalização por UE (média de diag(Hbar Hbar^H) = r).

% NOVO
A SNR é definida como $10\log_{10}(E_s/\sigma^2)$ \cite{zhao2025jsac}. Assim, só a razão entre $E_s$ e $\sigma^2$ importa. Por isso foi usado $E_s = 1$, com as constelações normalizadas para energia média unitária. Dessa forma, a potência do ruído térmico é $\sigma^2 = 10^{-\mathrm{SNR}/10}$.
% FONTE: L577 (definição da SNR).
% NOSSO (derivação): na eq. (3), multiplicar E_s e R pelo mesmo fator não muda W.
%   Como beta é proporcional a sigma^2 (IoT), R é proporcional a sigma^2. Então só
%   E_s/sigma^2 importa. Se a banca pedir, mostrar essa conta em uma linha.

% NOVO
A IoT é definida como $10\log_{10}((\beta + N_0)/\sigma^2)$ \cite{zhao2025jsac}, em que $N_0$ é a potência do ruído térmico, $\sigma^2$. O parâmetro $\beta$ é a potência de cada interferente, mas cada antena recebe a soma dos $r$ interferentes. Por isso, $\beta$ foi dividido por $r$. Assim, a razão entre a interferência e o ruído térmico em cada antena fica igual à IoT da configuração. Sem essa divisão, com IoT de 10~dB e $r = 8$, essa razão seria de cerca de 18,6~dB. Essa forma de medir a interferência, pela média por antena, também é usada por \citeonline{tomasin-interference}.
% FONTE: L577 (definição da IoT; N_0 não aparece em outro lugar do artigo).
% TODO FONTE: N_0 = sigma^2 vem do artigo anterior do mesmo grupo (Zhao et al.,
%   GLOBECOM 2021, L207: SNR medida como Es/N0). Esse artigo não está no .bib.
% NOSSO (conta): com a normalização por UE, a média de diag(Hbar Hbar^H) é r.
%   Sem dividir: 10*log10((10^1 - 1)*8 + 1) = 10*log10(73) = 18,6 dB.
% FONTE: Tomasin et al., eq. (47), linha 515 de Estimation*.md: ROT = (tr R / N)/sigma^2 - 1.
%   CUIDADO: o ROT deles é interferência/ruído; a IoT do artigo é
%   (interferência + ruído)/ruído. IoT linear = 1 + ROT. Por isso o texto diz só
%   "medir pela média por antena", não "IoT = ROT".

% NOVO
As amostras usadas para estimar $\hat{\mathbf{R}}$ vêm dos REs de piloto \cite{zhao2025jsac}. Por isso, o ruído dos símbolos de dados é sorteado de forma independente dessas amostras. Assim, o equalizador é testado com um ruído que não foi usado na estimação.
% FONTE: L123 ("noise samples in N pilot resource elements").
% NOSSO: a escolha (params.M5_Noise.data_noise_independent, status A/2B).
%   (Zhao et al. não tratam disso; não citar como deles.)


% ---------------------------------------------------------------------
\subsection{Receptor}\label{subsec:setup_receptor}
% ---------------------------------------------------------------------

% NOVO
A Tabela~\ref{tab:param_receptor} mostra os parâmetros da divisão das antenas, da estimação da covariância e dos equalizadores.

\begin{table}[htbp]
\caption{Parâmetros do receptor: particionamento, estimação da covariância e equalização.}
\label{tab:param_receptor}
\centering
\footnotesize
\renewcommand{\arraystretch}{1.2}
\begin{tabular}{|p{2.2cm}|p{3.6cm}|p{4.6cm}|p{3.4cm}|}
\hline
\raggedright Origem & \raggedright Parâmetro & \raggedright Nome no código & \raggedright Valor \tabularnewline
\hline
\raggedright \cite{zhao2025jsac} & \raggedright Antenas por \textit{cluster} & \raggedright \texttt{M3\_\allowbreak{}Clustering.\allowbreak equal\_\allowbreak{}cluster\_\allowbreak{}size} & \raggedright $M_c = M/C$ \tabularnewline
\cline{2-4}
 & \raggedright Amostras de ruído para estimar $\hat{\mathbf{R}}$ & \raggedright \texttt{M5\_\allowbreak{}Noise.\allowbreak n\_\allowbreak{}noise\_\allowbreak{}samples} & \raggedright \sout{$N = 192$} \textcolor{red}{No máximo $N = 192$} \tabularnewline
\cline{2-4}
 & \raggedright Estimativa da covariância & \raggedright \texttt{M6\_\allowbreak{}Covariance.\allowbreak estimator} & \raggedright Média amostral das $N$ amostras \tabularnewline
\cline{2-4}
 & \raggedright Conhecimento do canal na BS & \raggedright \texttt{M7\_\allowbreak{}Equalizers.\allowbreak perfect\_\allowbreak{}csi} & \raggedright $\mathbf{H}$ perfeito \tabularnewline
\cline{2-4}
 & \raggedright Equalizadores comparados & \raggedright \texttt{M7\_\allowbreak{}Equalizers.\allowbreak curves} & \raggedright LMMSE, ZF e BCD-MMSE \tabularnewline
\cline{2-4}
 & \raggedright Iterações do BCD-MMSE & \raggedright \texttt{M7\_\allowbreak{}Equalizers.\allowbreak bcd\_\allowbreak{}iterations} & \raggedright 1 \tabularnewline
\cline{2-4}
 & \raggedright Ordem de atualização do BCD & \raggedright \texttt{M7\_\allowbreak{}Equalizers.\allowbreak bcd\_\allowbreak{}update\_\allowbreak{}rule} & \raggedright Sequencial, $c = 1, \dots, C$ \tabularnewline
\cline{2-4}
 & \raggedright Cálculo do equalizador $\mathbf{W}$ & \raggedright \texttt{M1\_\allowbreak{}Driver.\allowbreak recompute\_\allowbreak{}W\_\allowbreak{}per\_\allowbreak{}realization} & \raggedright Uma vez por realização e por ponto de SNR \tabularnewline
\hline
\raggedright Este trabalho & \raggedright Divisão das antenas em \textit{clusters} & \raggedright \texttt{M3\_\allowbreak{}Clustering.\allowbreak mapping} & \raggedright Contígua \tabularnewline
\cline{2-4}
 & \raggedright \textcolor{red}{$N$ em cada condição de teste} & \raggedright \texttt{M5\_\allowbreak{}Noise.\allowbreak n\_\allowbreak{}noise\_\allowbreak{}samples\_\allowbreak{}by\_\allowbreak{}condition} & \raggedright \textcolor{red}{12, 48, 144 ou 192 (Tabela~\ref{tab:n-maximo})} \tabularnewline
\cline{2-4}
 & \raggedright Sorteio das amostras de ruído & \raggedright \texttt{M5\_\allowbreak{}Noise.\allowbreak resample\_\allowbreak{}noise\_\allowbreak{}per\_\allowbreak{}realization} & \raggedright Novas a cada realização e ponto de SNR \tabularnewline
\cline{2-4}
 & \raggedright Condicionamento de $\hat{\mathbf{R}}$ & \raggedright \texttt{M6\_\allowbreak{}Covariance.\allowbreak shrinkage} & \raggedright \textit{Shrinkage} de Ledoit--Wolf \tabularnewline
\cline{2-4}
 & \raggedright Forma do ZF & \raggedright \texttt{M7\_\allowbreak{}Equalizers.\allowbreak zf\_\allowbreak{}form} & \raggedright $(\mathbf{H}^H\mathbf{H})^{-1}\mathbf{H}^H$ \tabularnewline
\hline
\end{tabular}
\legend{Fonte: \citeonline{zhao2025jsac} e a autora. $C$ e $K$ estão na Tabela~\ref{tab:paineis}. Nomes no código: \texttt{dbp\_sim/params.py}.} % TODO citar Ledoit e Wolf (2004)
\end{table}
% FONTES: M_c = M/C L577, L85; N L577, L640; covariância eq. (6) L119 e L125;
%   H perfeito L63 e nota 2 (L817); equalizadores L581 e Figura 8;
%   iterações e posto L587, L591; Gauss-Seidel L398, nota 6 (L450);
%   L = K L150, L184, L593; W reaproveitado L61, L561; ZF citado em L581 sem fórmula.
% CORREÇÃO em relação à tabela antiga: W e as amostras de ruído são refeitos a
%   cada ponto de SNR, não só a cada realização (m1_driver.run_panel).
% LEMBRETE: o Ledoit-Wolf também é aplicado no R_cc do BCD (m7_equalizers.bcd_mmse).
%   Com N = 192 o bloco M_c x M_c tem inversa; com N = 12 (HST) não tem, pois
%   N < M_c. O módulo alvo do constraints.csv inverte o R_cc sem ajuste.
% N por condição: params.M5_Noise.n_noise_samples_by_condition; run_figure8.py --conditions.

% NOVO
Cada \textit{cluster} fica com antenas vizinhas. O primeiro fica com as antenas 1 a $M_c$, o segundo com as antenas $M_c+1$ a $2M_c$, e assim por diante. Essa divisão segue a notação em blocos $\mathbf{H}=[\mathbf{H}_1^T,\dots,\mathbf{H}_C^T]^T$ de \citeonline{zhao2025jsac}.
% FONTE: L85 (divisão em C clusters); L96 (notação em blocos).
% Outra fonte possível: o código de referência DBP-master (DBP_detector_sim.m:305)
%   e o artigo dele (ref. [6] do artigo base, nota 3). Não está no .bib.

% NOVO
O equalizador $\mathbf{W}$ é calculado uma vez para cada realização de canal e cada ponto de SNR. Depois, ele é usado para todos os 480 símbolos. Uma mesma matriz de equalização pode ser usada para vários símbolos, enquanto o canal não muda \cite{zhao2025jsac}.

% NOVO (2026-10-14)
\textcolor{red}{O número $N$ de amostras de ruído depende da condição de teste. Ele é o número de REs de DM-RS de um PRB dentro de $T_c$, limitado a 192 (Tabela~\ref{tab:n-maximo}). Assim, cada configuração da Tabela~\ref{tab:paineis} é simulada com $N = 192$, 144, 48 e 12. A simulação não tem eixo de tempo e o canal não tem Doppler. Por isso, entre as condições de teste, só muda o valor de $N$.}
% FONTE: Tabela tab:n-maximo (dbp-ruido.tex); código: params.M5_Noise.n_noise_samples_by_condition,
%   run_figure8.py --conditions all. "Sem Doppler": nota do próprio params.py.
% FONTE: L61, L561 (reuso da matriz de equalização).
% Código: m1_driver.run_panel (W montado dentro do laço de SNR).

% AUTORA (9-results.tex:333), tom ajustado e "soma" trocada por "média ponderada"
% ATENÇÃO (não alterado): "soma à matriz um múltiplo da identidade" descreve a
%   carga diagonal fixa, a alternativa rejeitada. O Ledoit-Wolf faz uma média
%   ponderada: (1 - rho) R_hat + rho mu I (m5_noise.shrink). Decidir a redação.
O equalizador LMMSE usa a inversa de $\hat{\mathbf{R}}$ \cite{zhao2025jsac}. Com $M = 256$ antenas e $N = 192$ amostras, $\hat{\mathbf{R}}$ não tem inversa. \textcolor{red}{Com $N = 48$ ou $N = 12$, isso vale também para $M = 128$.} Por isso foi usado o \textit{shrinkage} de Ledoit--Wolf. Ele faz uma média ponderada entre $\hat{\mathbf{R}}$ e um múltiplo da identidade, e o peso dessa média é calculado a partir das próprias amostras. Assim não foi preciso escolher uma constante. % TODO citar Ledoit e Wolf (2004)

% NOVO
No BCD-MMSE, cada DU inverte a matriz $E_s\mathbf{H}_c\mathbf{H}_c^H + \hat{\mathbf{R}}_{cc}$, de tamanho $M_c \times M_c$ \cite{zhao2025jsac}. \sout{Nas configurações simuladas, $M_c \le 32$, um valor menor que $N = 192$. Mesmo assim, o simulador também aplica o \textit{shrinkage} ao bloco $\hat{\mathbf{R}}_{cc}$ antes dessa inversão.} \textcolor{red}{Nas configurações simuladas, $M_c \le 32$. Com $N = 192$, 144 ou 48, esse valor é menor que $N$. Nas condições HST, com $N = 12$, ele é maior que $N$, e o bloco $\hat{\mathbf{R}}_{cc}$ não tem inversa (Seção~\ref{subsec:ruido-colorido}). O simulador aplica o \textit{shrinkage} ao bloco $\hat{\mathbf{R}}_{cc}$ em todas as condições.}
% FONTE: eq. (30), L504; tamanhos na Tabela tab:paineis.
% Código: m7_equalizers.bcd_mmse -- den = Es*Hc*Hc^H + cond.covariance(Sc).
% DECIDIR ANTES DE MANTER: o módulo alvo do constraints.csv inverte o R_cc sem
%   ajuste. Duas saídas: (1) manter o shrinkage e dizer isso aqui; (2) rodar o BCD
%   com --shrinkage none, para os resultados mostrarem o que o hardware vai receber.
% SEM FONTE: dizer que o bloco "tem inversa" porque N > M_c precisa de referência
%   (ligar à seção "Quantidade de amostras" do capítulo 1). Por isso a frase só
%   compara os números.

% NOVO
O ZF centralizado é usado como referência \cite{zhao2025jsac}, na forma da Equação~\eqref{eq:zf_filter}.
% FONTE: L581 (ZF citado sem fórmula). eq:zf_filter está em 1-noise-and-constraints.tex.


% ---------------------------------------------------------------------
\subsection{Avaliação}\label{subsec:setup_avaliacao}
% ---------------------------------------------------------------------

% NOVO
A Tabela~\ref{tab:param_avaliacao} mostra como os resultados são medidos.

\begin{table}[htbp]
\caption{Parâmetros da avaliação.}
\label{tab:param_avaliacao}
\centering
\footnotesize
\renewcommand{\arraystretch}{1.2}
\begin{tabular}{|p{2.2cm}|p{3.6cm}|p{4.6cm}|p{3.4cm}|}
\hline
\raggedright Origem & \raggedright Parâmetro & \raggedright Nome no código & \raggedright Valor \tabularnewline
\hline
\raggedright \cite{zhao2025jsac} & \raggedright Realizações de canal & \raggedright \texttt{M1\_\allowbreak{}Driver.\allowbreak n\_\allowbreak{}realizations} & \raggedright 100 \tabularnewline
\cline{2-4}
 & \raggedright Símbolos por realização & \raggedright \texttt{M1\_\allowbreak{}Driver.\allowbreak n\_\allowbreak{}symbols} & \raggedright 480 \tabularnewline
\cline{2-4}
 & \raggedright Detecção & \raggedright \texttt{M8\_\allowbreak{}Metric.\allowbreak detection} & \raggedright Ponto mais próximo da constelação \tabularnewline
\cline{2-4}
 & \raggedright Métrica & \raggedright \texttt{M8\_\allowbreak{}Metric.\allowbreak metric} & \raggedright Taxa de erro de símbolo (SER) \tabularnewline
\cline{2-4}
 & \raggedright Faixa de SNR & \raggedright \texttt{M1\_\allowbreak{}Driver.\allowbreak snr\_\allowbreak{}db\_\allowbreak{}grid} & \raggedright $-10$ a $10$ dB, passo de 2 dB \tabularnewline
\hline
\raggedright Este trabalho & \raggedright Remoção de viés antes da detecção & \raggedright \texttt{M8\_\allowbreak{}Metric.\allowbreak bias\_\allowbreak{}removal} & \raggedright Não \tabularnewline
\cline{2-4}
 & \raggedright Semente aleatória & \raggedright \texttt{M1\_\allowbreak{}Driver.\allowbreak rng\_\allowbreak{}seed} & \raggedright 0 \tabularnewline
\hline
\end{tabular}
\legend{Fonte: \citeonline{zhao2025jsac} e a autora. Nomes no código: \texttt{dbp\_sim/params.py}.}
\end{table}

% NOVO
Cada ponto das curvas é a média da SER sobre 100 realizações de canal, com 480 símbolos em cada uma. A SNR varia de $-10$ a $10$~dB, com passo de 2~dB, como em \citeonline{zhao2025jsac}.
% FONTE: L577 (100 x 480); faixa de SNR lida dos marcadores das Figs. 6 a 8
%   (o passo não está escrito no artigo; ver params.M1_Driver.snr_db_grid).

% NOVO
Os símbolos são detectados pelo ponto mais próximo da constelação \cite{zhao2025jsac}. Antes da detecção, a estimativa de cada UE não passa por nenhuma correção de ganho. A semente do gerador aleatório foi fixada em 0, para que os resultados possam ser repetidos.
% FONTE: L79 (detecção). "Não fala dessa correção": busca registrada em
%   params.M8_Metric.bias_removal. Código: m8_metric.symbol_error_rate.
